\sect{The transition values}
\label{sec:endpoint}

At a Hermite zero the optimal exponent lies halfway between the exponents on
the two neighboring open intervals. For each fixed $d\geq 0$, with $a=h_{d+1}$,
we have
\begin{equation}
  b_k(a)=\Theta_d(k^{-d-1}).                                      \label{eq:endpoint}
\end{equation}
This includes $a=h_1=0$. We give the details because the support condition
at equality makes direct rounding of a Gaussian quadrature rule delicate.

\paragraph{Lower bound.}
Put $q(x)=H_{d+1}(x)/(x-a)$ and $F_k(x)=(x+\sqrt{k})(x-a)q(x)^2$.
On the permitted part of the cube $-\sqrt{k}\le x\le a$, we have $F_k(x)\le0$.
Write $Z$ for a standard Gaussian. Orthogonality to $q$ gives
$\mathbb E[(Z-a)q(Z)^2]=0$, and consequently
\[
  \mathbb E[Z(Z-a)q(Z)^2]=\mathbb E[H_{d+1}(Z)^2]=(d+1)!.
\]
For every fixed polynomial the uniform-cube expectation at $X$ equals its
Gaussian expectation up to $O_d(k^{-1})$ (with constants depending on that
polynomial). Indeed the moments of odd order vanish, and in each even moment
the pairings of distinct indices give the Gaussian main term with error
$O_d(k^{-1})$. Hence $\mathbb E_{\rm unif}F_k(X)=(d+1)!+O_d(k^{-1/2})$.
The sum of the absolute Walsh coefficients of $X^j$ is at most $k^{j/2}$,
by expanding the $j$-fold product. It follows term by term that the sum for
$F_k(X)$ is $O_d(k^{d+1})$. Comparing the uniform expectation with the
expectation for any supported law proves the lower bound in
\eqref{eq:endpoint}.

\paragraph{An interior quadrature rule.}
Let $m=d+1$ and write $z_1<\cdots<z_m=a$ for the zeros of $H_m$.
Let $J(t)$ be the $m$-by-$m$ tridiagonal Hermite matrix with off-diagonal
entries $\sqrt{1},\ldots,\sqrt{m-1}$ and diagonal entries zero except for
the last, which is $t$. Its spectral distribution at the first coordinate
has positive weights at its $m$ simple eigenvalues $z_i(t)$. The eigenvector
recurrence shows that both endpoint coordinates are nonzero, and
$z_i'(0)$ is the square of the last coordinate of a normalized eigenvector.
In particular $z_i'(0)>0$. Spectral weights are continuous near $0$.
Every one of these distributions integrates the Gaussian moments of orders
$0$ through $2d$. Indeed the moments are $(J(t)^j)_{11}$, and a walk from
the first coordinate using the final diagonal entry and returning needs
at least $2d+1$ steps. The rule at $t=0$ is Gaussian quadrature.

Choose a constant $L>12/z_m'(0)$ and take $t=-L/\sqrt{k}$.
For all sufficiently large $k$, even the largest node satisfies
$z_m(t)<a-10/\sqrt{k}$. Each node is within $L/\sqrt{k}$ of its
corresponding zero $z_i$, by the operator norm bound on $J(t)-J(0)$.
The nodes remain separated, and their weights stay bounded below by a
positive constant depending on $d$.

\clearpage
\paragraph{Exact moment matching after rounding.}
Write $w$ for the number of minus signs, so $X=(k-2w)/\sqrt{k}$.
Fix one of the two parities of $w$. For each node choose two attainable
values of $X$ of this parity, one on either side at distance between
$4/\sqrt{k}$ and $8/\sqrt{k}$. This is possible because the mesh is
$4/\sqrt{k}$, and all these values still lie below $a$. Split each
original weight between its two neighbors by linear interpolation,
preserving the original node's mean. Both fractions are at least $1/3$.
For every fixed monomial of degree at most $2d+1$, Taylor's theorem on a
fixed compact interval gives a total interpolation error $O_d(k^{-1})$.

At the $2m=2d+2$ grid nodes prescribe the exact uniform-cube moments through
$2d$, and prescribe the original $J(t)$ moment for degree $2d+1$. The
split weights miss each of these targets by $O_d(k^{-1})$. Through $2d$
the original rule has the Gaussian moments, whose discrepancies from the
cube moments are $O_d(k^{-1})$. Correct all the weights by solving the
square Vandermonde system for these errors. The inverse has entries
$O_d(\sqrt{k})$. Indeed the coefficients of the Lagrange basis polynomials
have this bound, since each denominator has one factor of order
$k^{-1/2}$ and all the other factors stay bounded away from zero, while the
numerator coefficients stay bounded. Thus each weight changes by
$O_d(k^{-1/2})$ and remains positive. The zeroth moment keeps their sum one.
Carry out this construction on each parity, spread each mass uniformly over
its layer, and average the two probability laws. In each parity the moments
through $2d$ match the cube exactly, and all the supporting values lie
within $O_d(k^{-1/2})$ of a zero of $H_{d+1}$.

\paragraph{Bias on every level.}
Let $K_j(w)$ be the sum of the Walsh characters of size $j$ on the layer
$w$. Then $P_{j,k}(x)=j!k^{-j/2}K_j(w)$ is monic in $x=X$. For fixed $j$
it stays bounded on compact intervals, as follows from $P_{0,k}=1$,
$P_{1,k}=x$, and
\[
 P_{j+1,k}=xP_{j,k}-j\bigl(1-(j-1)/k\bigr)P_{j-1,k}.
\]
For $1\le j\le 2d$ its expectation on each parity is zero by exact moment
matching. For $j=2d+1$, the polynomial $P_{j,k}-H_{d+1}H_d$ has degree at
most $2d$ and uniform-cube mean zero by oddness. On the supporting layers
$H_{d+1}H_d=O_d(k^{-1/2})$. Thus the expectation of $P_{2d+1,k}$ on each
parity is $O_d(k^{-1/2})$. Dividing the expected $K_j$ by $\binom{k}{j}$
gives the bias of each character of that size. The bound for $j=2d+1$ is
therefore $O_d(k^{-d-1})$, and the same bound holds pointwise on the
support for $j=2d+2$ by the fixed-degree recurrence. Reflection
$K_{k-j}(w)=(-1)^wK_j(w)$ controls the corresponding upper levels, with
the top level zero by equal parity masses.

For $2d+3\le j\le k-2d-3$, orthogonality of the layer characters gives
\[
 \frac{|K_j(w)|}{\binom{k}{j}}
  \le \left(\frac{2^k}{\binom{k}{w}\binom{k}{j}}\right)^{1/2}
  = O_d(k^{-d-5/4}).
\]
Here $|k-2w|=O_d(\sqrt{k})$, so $2^k/\binom{k}{w}=O_d(\sqrt{k})$ by
Stirling's formula. This proves the upper bound in \eqref{eq:endpoint}.
